\documentclass[a4paper,11pt,oneside]{memoir} 	% Openright aabner kapitler paa hoejresider (openany = vilkaarlig/begge)fleqn,
%--------------------------------------------------------------------------
%%%% PAKKER %%%%
%--------------------------------------------------------------------------
% ¤¤ Oversaettelse og tegnsaetning ¤¤ %
\usepackage[utf8]{inputenc}					% Input-indkodning af tegnsaet, dvs. input fra keyboard, tegnoversigt eller andet (UTF8 = Unicode)
\usepackage[T1]{fontenc}					% Output-indkodning af tegnsaet, dvs. printede fonte og tegn (T1 = Type 1 font med support for de fleste europaeiske sprog)
\usepackage[danish]{babel}					% Sproglig fremstilling af elementer (figur vs. figure, litteratur vs. bibliography osv.)
\usepackage{ragged2e,anyfontsize}			% Justering af elementer
%--------------------------------------------------------------------------		
% ¤¤ Figurer og tabeller (floats) ¤¤ %
\usepackage{rotating, graphicx} 			% Inkludering af eksterne billeder (JPG, PNG, PDF)
\usepackage{multirow}                		% Fletning af raekker og kolonner (\multicolumn og \multirow)
\usepackage{colortbl} 						% Farver i tabeller (fx \columncolor, \rowcolor og \cellcolor)
\usepackage[dvipsnames]{xcolor}				% Definer farver med \definecolor. Se mere: http://en.wikibooks.org/wiki/LaTeX/Colors
\usepackage{flafter}						% Soerger for, at floats ikke optraeder i teksten foer deres reference
\usepackage{float}							% Muliggoer eksakt placering af floats, fx \begin{figure}[H]
\let\newfloat\relax 						% Justering mellem float-pakken og memoir
%\usepackage{eso-pic}						% Tilfoej billedekommandoer paa hver side
%\usepackage{wrapfig}						% Indsaettelse af figurer omsvoebt af tekst 
%\usepackage{multicol}         	        	% Muliggoer tekst i spalter
%\usepackage{rotating}						% Rotation af tekst med \begin{sideways}...\end{sideways}
%--------------------------------------------------------------------------
% ¤¤ Matematik mm. ¤¤
\usepackage{amsmath,amssymb,stmaryrd} 		% Avancerede matematik-udvidelser
\usepackage{mathtools}						% Andre matematik- og tegnudvidelser
\usepackage{textcomp}                 		% Symbol-udvidelser (fx promille-tegn med \textperthousand)
\usepackage[version=3]{mhchem} 			% Kemi-pakke til flot og let notation af formler, fx \ce{Fe2O3}
%\usepackage{rsphrase}						% Kemi-pakke til RS-saetninger, fx \rsphrase{R1}
%--------------------------------------------------------------------------
% ¤¤ Referencer og kilder ¤¤ %
\usepackage[danish]{varioref}				% Muliggoer bl.a. krydshenvisninger med sidetal (\vref)
\usepackage{natbib}							% Udvidelse med naturvidenskabelige citationsmodeller, herunder Harvard-modellen
%\usepackage{xr}							% Referencer til eksternt dokument med \externaldocument{<NAVN>}
%\usepackage{glossaries}					% Terminologi- eller symbolliste (se mere i Lars Madsens Latex-bog)
%--------------------------------------------------------------------------
% ¤¤ Misc. ¤¤ %
\usepackage{listings}						% Placer kildekode i dokumentet med \begin{lstlisting}...\end{lstlisting}
\usepackage{lipsum}							% Dummy tekst med fx \lipsum[2]
\usepackage[shortlabels]{enumitem}			% Muliggoer enkelt konfiguration af lister (se \setlist nedenfor)
\usepackage{pdfpages}						% Goer det muligt at inkludere pdf-dokumenter med kommandoen \includepdf[pages={x-y}]{fil.pdf}	
%\pdfoptionpdfminorversion			% Muliggoer inkludering af pdf-dokumenter af version 1.6 og hoejere
\pretolerance=2500 							% Justering af afstand mellem ord (hoejt tal, mindre orddeling og mere luft mellem ord)

% Kommentarer og rettelser med \fxnote. Med 'final' i stedet for 'draft' udloeser hver note en error i den faerdige rapport.
\usepackage[footnote,draft,danish,silent,nomargin]{fixme}		
%--------------------------------------------------------------------------
%%%% BRUGERDEFINEREDE INDSTILLINGER %%%%
%--------------------------------------------------------------------------
% ¤¤ Marginer ¤¤ %
\setlrmarginsandblock{3.5cm}{2.5cm}{*}		% \setlrmarginsandblock{Indbinding}{Kant}{Ratio}
\setulmarginsandblock{2.5cm}{3.0cm}{*}		% \setulmarginsandblock{Top}{Bund}{Ratio}
\checkandfixthelayout 						% Oversaetter vaerdier til brug for andre pakker
%--------------------------------------------------------------------------
%	¤¤ Afsnitsformatering ¤¤ %
\setlength{\parindent}{0mm}           		% Stoerrelse af indryk
\setlength{\parskip}{3mm}          			% Afstand mellem afsnit ved brug af double Enter
\linespread{1,1}							% Linjeafstand
%--------------------------------------------------------------------------
% ¤¤ Litteraturlisten ¤¤ %
% \bibpunct[,]{[}{]}{;}{a}{,}{,} 				% Definerer parametre ved Harvard-henvisning (bl.a. parantestype og seperatortegn)
% \bibliographystyle{bibtex/plain}			% Udseende af litteraturlisten (Harvard-metoden - skift til fx 'plain' for tal)
%--------------------------------------------------------------------------
% ¤¤ Lister ¤¤ %
\setlist{
  topsep=0pt,								% Vertikal afstand mellem tekst og listen
  itemsep=-1ex,								% Vertikal afstand mellem items
} 
%--------------------------------------------------------------------------
% ¤¤ Visuelle referencer ¤¤ %
\usepackage[colorlinks]{hyperref}			% Danner klikbare referencer (hyperlinks) i dokumentet
\hypersetup{colorlinks = true,				% Opsaetning af farvede hyperlinks (interne links, citeringer og URL)
    linkcolor = black,
    citecolor = black,
    urlcolor = black
}
%--------------------------------------------------------------------------
% ¤¤ Opsaetning af figur- og tabeltekst ¤¤ %
\captionnamefont{\small\bfseries\itshape}	% Opsaetning af tekstdelen ('Figur' eller 'Tabel')
\captiontitlefont{\small}					% Opsaetning af nummerering
\captiondelim{. }							% Seperator mellem nummerering og figurtekst
\captionstyle{\centering}					% Justering/placering af figurteksten (centreret = \centering, venstrejusteret = \raggedright)
\captionwidth{\linewidth}					% Bredden af figurteksten
\hangcaption								% Venstrejusterer fler-linjers figurtekst under hinanden
\setlength{\belowcaptionskip}{0pt}			% Afstand under figurteksten
%--------------------------------------------------------------------------		
% ¤¤ Opsaetning af listings ¤¤ %
\definecolor{commentGreen}{RGB}{34,139,24}
\definecolor{stringPurple}{RGB}{208,76,239}
%--------------------------------------------------------------------------
\lstset{language=Matlab,					% Sprog
	basicstyle=\ttfamily\scriptsize,		% Opsaetning af teksten
	keywords={for,if,while,else,elseif,		% Noegleord at fremhaeve
			  end,break,return,case,
			  switch,function},
	keywordstyle=\color{blue},				% Opsaetning af noegleord
	commentstyle=\color{commentGreen},		% Opsaetning af kommentarer
	stringstyle=\color{stringPurple},		% Opsaetning af strenge
	showstringspaces=false,					% Mellemrum i strenge enten vist eller blanke
	numbers=left, numberstyle=\tiny,		% Linjenumre
	extendedchars=true, 					% Tillader specielle karakterer
	columns=flexible,						% Kolonnejustering
	breaklines, breakatwhitespace=true,		% Bryd lange linjer
}
%--------------------------------------------------------------------------
% ¤¤ Navngivning ¤¤ %
\addto\captionsdanish{
	\renewcommand\contentsname{Indholdsfortegnelse}			% Skriver 'Indholdsfortegnelse' i stedet for 'Indhold'
	\renewcommand\appendixname{Appendiks}					% Skriver 'Appendiks' i stedet for 'Appendix'
	\renewcommand\appendixpagename{Appendiks}
	\renewcommand\appendixtocname{Appendiks}
	\renewcommand\cftchaptername{\chaptername~}				% Skriver 'Kapitel' foran kapitlerne i indholdsfortegnelsen
	\renewcommand\cftappendixname{\appendixname~}			% Skriver 'Appendiks' foran appendiks i indholdsfortegnelsen
}
%--------------------------------------------------------------------------
% ¤¤ Kapiteludssende ¤¤ %
\definecolor{numbercolor}{gray}{0.7}		% Definerer en farve til brug til kapiteludseende
\newif\ifchapternonum

\makechapterstyle{jenor}{					% Definerer kapiteludseende frem til ...
  \renewcommand\beforechapskip{0pt}
  \renewcommand\printchaptername{}
  \renewcommand\printchapternum{}
  \renewcommand\printchapternonum{\chapternonumtrue}
  \renewcommand\chaptitlefont{\fontfamily{pbk}\fontseries{db}\fontshape{n}\fontsize{25}{35}\selectfont\raggedleft}
  \renewcommand\chapnumfont{\fontfamily{pbk}\fontseries{m}\fontshape{n}\fontsize{1in}{0in}\selectfont\color{numbercolor}}
  \renewcommand\printchaptertitle[1]{%
    \noindent
    \ifchapternonum
    \begin{tabularx}{\textwidth}{X}
    {\let\\\newline\chaptitlefont ##1\par} 
    \end{tabularx}
    \par\vskip-2.5mm\hrule
    \else
    \begin{tabularx}{\textwidth}{Xl}
    {\parbox[b]{\linewidth}{\chaptitlefont ##1}} & \raisebox{-15pt}{\chapnumfont \thechapter}
    \end{tabularx}
    \par\vskip2mm\hrule
    \fi
  }
}											% ... her

\chapterstyle{jenor}						% Valg af kapiteludseende - Google 'memoir chapter styles' for alternativer
%--------------------------------------------------------------------------
% ¤¤ Sidehoved/sidefod ¤¤ %
\makepagestyle{Uni}							% Definerer sidehoved og sidefod udseende frem til ...
\makepsmarks{Uni}{%
	\createmark{chapter}{left}{shownumber}{}{. \ }
	\createmark{section}{right}{shownumber}{}{. \ }
	\createplainmark{toc}{both}{\contentsname}
	\createplainmark{lof}{both}{\listfigurename}
	\createplainmark{lot}{both}{\listtablename}
	\createplainmark{bib}{both}{\bibname}
	\createplainmark{index}{both}{\indexname}
	\createplainmark{glossary}{both}{\glossaryname}
}
\nouppercaseheads											% Ingen Caps oenskes

\makeevenhead{Uni}{}{}{\leftmark}				% Lige siders sidehoved (\makeevenhead{Navn}{Venstre}{center}{Hoejre})
\makeoddhead{Uni}{\rightmark}{}{}	% Ulige siders sidehoved (\makeoddhead{Navn}{Venstre}{Center}{Hoejre})
%\makeevenfoot{Uni}{\thepage}{}{}							% Lige siders sidefod (\makeevenfoot{Navn}{Venstre}{Center}{Hoejre})
%\makeoddfoot{Uni}{}{}{\thepage}								% Ulige siders sidefod (\makeoddfoot{Navn}{Venstre}{Center}{Hoejre})
\makeheadrule{Uni}{\textwidth}{0.5pt}						% Tilfoejer en streg under sidehovedets indhold
\makefootrule{Uni}{\textwidth}{0.5pt}{1mm}					% Tilfoejer en streg under sidefodens indhold

\copypagestyle{Unichap}{Uni}								% Der dannes en ny style til kapitelsider
\makeoddhead{Unichap}{}{}{}									% Sidehoved defineres som blank på kapitelsider
\makeevenhead{Unichap}{}{}{}
\makeheadrule{Unichap}{\textwidth}{0pt}
\aliaspagestyle{chapter}{Unichap}							% Den ny style vaelges til at gaelde for chapters
															% ... her
\pagestyle{uni}												% Valg af sidehoved og sidefod (benyt 'plain' for ingen sidehoved/fod). 'uni' er normal
%--------------------------------------------------------------------------
%%%% EGNE KOMMANDOER %%%%

% ¤¤ Billede hack ¤¤ %										% Indsaet figurer nemt med \figur{Stoerrelse}{Fil}{Figurtekst}{Label}
\newcommand{\figur}[4]{
		\begin{figure}[H] \centering
			\includegraphics[width=#1\textwidth]{billeder/#2}
			\caption{#3}
			\label{#4}
		\end{figure} 
}
%--------------------------------------------------------------------------
% ¤¤ Specielle tegn ¤¤ %
\newcommand{\dec}{^{\circ}}									% '\dec' returnerer et gradtegn (husk $$ udenfor aligns)
\newcommand{\decC}{^{\circ}\text{C}}						% '\decC' returnerer et gradtegn + 'C' (husk $$ udenfor aligns)
\newcommand{\m}{\cdot}										% '\m' returnerer et gangetegn
%--------------------------------------------------------------------------
%%%% ORDDELING %%%%
%--------------------------------------------------------------------------
\hyphenation{In-te-res-se e-le-ment}
%%%%%%%%%%%%%%%%%%%%%%%%%%%%%%%%%%%%%%%%%%%%%%%%%%%%%%%%%%%%%%%%%%%%%%%%%%%%%%
%%%%%%%%%%%%%%%%%%%%%%%%%%%%% -MINE ÆNDRINGER- %%%%%%%%%%%%%%%%%%%%%%%%%%%%%%%
%%%%%%%%%%%%%%%%%%%%%%%%%%%%% -MINE ÆNDRINGER- %%%%%%%%%%%%%%%%%%%%%%%%%%%%%%%
%%%%%%%%%%%%%%%%%%%%%%%%%%%%% -MINE ÆNDRINGER- %%%%%%%%%%%%%%%%%%%%%%%%%%%%%%%
%%%%%%%%%%%%%%%%%%%%%%%%%%%%% -MINE ÆNDRINGER- %%%%%%%%%%%%%%%%%%%%%%%%%%%%%%%
%%%%%%%%%%%%%%%%%%%%%%%%%%%%% -MINE ÆNDRINGER- %%%%%%%%%%%%%%%%%%%%%%%%%%%%%%%
%%%%%%%%%%%%%%%%%%%%%%%%%%%%%%%%%%%%%%%%%%%%%%%%%%%%%%%%%%%%%%%%%%%%%%%%%%%%%%
%--------------------------------------------------------------------------\textbf{}
% EKSTRA PAKKER:
%-------------------------------------------------------------------------- 
\usepackage{siunitx}    % SI-enheder: \si{enhed} og \SI{tal}{enhed} - BARAKAT
\sisetup{output-decimal-marker = {,}} % Opsaetning af \SI og decimalseparator
\usepackage{wrapfig}    % Tekst rundt om billeder                   - BARAKAT
\usepackage{caption}    % Billed- og tabeltekst                     - BARAKAT
\usepackage{titlesec}   % Styrer plads før/efter overskriver
\usepackage{blindtext}  % Genererer ligegyldigt brødtekst
                        % Bruges som fx: \blindtext[5]
\usepackage{natbib}     % Tillader Harvard som litteraturstandard
\usepackage[left=2cm, right=2cm, top=2.5cm, headsep=1cm]{geometry}     % Standard margin (word)
\usepackage{lastpage}
%--------------------------------------------------------------------------
% FORMATERING AF AFSTANDE FØR OG EFTER OVERSKRIFTER AF VILKÅRLIG DYBDE:
%-------------------------------------------------------------------------- 
\setcounter{tocdepth}{5}        % Styrer grad af underafsnit i t.o.c.   - BARAKAT
\settocdepth{subparagraph}      % Styrer grad af underafsnit i t.o.c.   - BARAKAT
\setsecnumdepth{subparagraph}   % Underafsnit af 5. grad.               - BARAKAT
\setcounter{secnumdepth}{5}     % Tal på Underafsnit af 5. grad         - BARAKAT
\renewcommand{\thesection}{\arabic{section}} % Afsnitnummering starter fra 1
\renewcommand{\thesection}{\thechapter.\arabic{section}} % Kapitelnummer medtages i titelnummerering i dokument og T.o.C.
%--------------------------------------------------------------------------
% FORMATERING AF AFSTANDE FØR OG EFTER OVERSKRIFTER AF VILKÅRLIG DYBDE:
%-------------------------------------------------------------------------- 
%\titlespacing{DYBDE}{<OVER> FAKTISK + -}{<UNDER> FAKTISK + -}
\titlespacing\section{0pt}{10pt plus 1pt minus 1pt}{-10pt plus 2pt minus 2pt}
\titlespacing\subsection{-300pt}{15pt plus 1pt minus 1pt}{-10pt plus 2pt minus 2pt}
\titlespacing\subsubsection{0pt}{15pt plus 1pt minus 1pt}{-10pt plus 2pt minus 2pt}
\titlespacing\paragraph{0pt}{15pt plus 1pt minus 1pt}{0pt plus 2pt minus 2pt}
%------------------------------------------------------------------
%--------------- Formatering af sidehoved- og sidefod -------------
%------------------------------------------------------------------
    \usepackage{fancyhdr}
    \pagestyle{fancy}
    \fancyhead[C]{\nouppercase{M4PRJ4 \\ Energitekniske systemer}}
    \fancyhead[L]{\nouppercase{IHA Maskinteknik, Gruppe 3 \\ 18-12-2018}}
    \fancyhead[R]{\nouppercase{\includegraphics[width=0.2\textwidth]{Billeder/AUlogo.JPG}}}
    \renewcommand\headrulewidth{1pt}
    \makeatletter
    \def\headrule{{\if@fancyplain\let\headrulewidth\plainheadrulewidth\fi
    \hrule\@height\headrulewidth\@width\headwidth
    \vskip 2pt                              % 2pt between lines
    \hrule\@height.5pt\@width\headwidth     % lower line with .5pt line width
    \vskip-\headrulewidth
    \vskip-1pt}}
    \makeatother
    \renewcommand{\footrulewidth}{1pt}
    
    %\fancypagestyle{plain}{% % <-- this is new
    %\fancyhf{} 
    %\fancyfoot[C]{Side \thepage\ af \pageref{LastPage}}
    %\renewcommand{\hrulewidth}{0pt}}
    
    \usepackage{titletoc}
    \usepackage{tocloft}
    
    %\fancyfoot[C]{Side \thepage\ af \pageref{LastPage}}
    %\usepackage{titlesec}
    %\assignpagestyle{\chapter}{plain}
%-------------------------------------------------------------------